\documentclass[12pt,a4paper,fontset=fandol]{ctexart}

\usepackage[margin=2.4cm]{geometry}
\usepackage{booktabs}
\usepackage{graphicx}
\graphicspath{{/Users/qiu/Desktop/cvfinal_full_submission/exports/report_assets/}{exports/report_assets/}{../../exports/report_assets/}}
\usepackage{hyperref}
\usepackage{amsmath}
\usepackage{float}

\title{计算机视觉期末作业：2DGS/AIGC 场景融合与 LeRobot ACT 泛化}
\author{阎丞麟 \quad 23307110060 \\
 潘艺赫\quad 23307110090}
\date{2026 年 6 月}

\begin{document}
\maketitle

\noindent GitHub 仓库：\url{TODO_PUBLIC_REPOSITORY_URL}

\noindent 模型权重与视频下载：\url{https://drive.google.com/drive/u/0/folders/1wO51nD2L4Q1x6lwkvZTFGSNTZrD8LXnQ}
\section{任务背景与数据集}
本文完成两个任务：基于 2DGS 与 AIGC 的多源 3D 资产生成和真实场景融合，以及基于 LeRobot ACT 的 CALVIN 跨环境泛化实验。

\textbf{CALVIN 数据集：} CALVIN (Composing Actions from Language and Vision) 是一个桌面机器人操控仿真环境，包含 4 种不同视觉布局的环境 A/B/C/D，每环境包含 34 个操控任务（如推方块、开门、抓取等）。数据使用 HuggingFace 上的 \texttt{xiaoma26/calvin-lerobot} 格式，经预处理后用于 LeRobot 框架训练。各 split 的划分如下：\texttt{splitA}、\texttt{splitB}、\texttt{splitC} 各含约 3,000 个 episode，用于训练；\texttt{splitD} 含约 150 个 episode，仅用于 zero-shot 测试。观测空间包括 200×200 的前置摄像头图像、84×84 的腕部相机图像以及 15 维本体状态（关节位置与夹爪宽度），动作空间为 7 维末端执行器控制指令。

\textbf{Mip-NeRF 360 数据集：} 背景场景选用 \texttt{counter} 场景，该场景包含约 200 张高分辨率照片，覆盖柜台区域的多个视角，用于 2DGS 重建作为统一环境背景。

\section{题目一：2DGS 与 AIGC 场景融合}
\subsection{方法原理}

\textbf{2D Gaussian Splatting (2DGS)：} 2DGS 是 3D Gaussian Splatting 的改进，将 3D 场景表示为大量带法线方向的 2D 高斯面片（surfels），结合可微光栅化进行渲染优化。给定一组多视角图像和对应的相机位姿（由 COLMAP 估计），2DGS 通过梯度下降优化各高斯面片的位置、颜色、不透明度和朝向，从随机初始化逐步收敛到与输入图像一致的 3D 重建。

\textbf{Threestudio / DreamFusion：} DreamFusion 利用预训练的 2D 文本-图像扩散模型（如 Imagen 或 Stable Diffusion）作为先验，通过 Score Distillation Sampling (SDS) Loss 指导 3D 表示（如 NeRF 或 Mesh）的优化。SDS 从当前渲染视角生成带噪声的图像，计算扩散模型估计的噪声残差，再反向传播到 3D 参数，实现"从文本到 3D"的零样本生成。

\textbf{Magic123：} Magic123 结合了 SDS 先验与参考图像的监督信号，在单张输入图像的约束下生成完整 3D 模型。其优化框架包含两个阶段：coarse 阶段使用 NeRF 表示进行粗略重建，fine 阶段通过 Mesh 表示精细化纹理。本实验仅使用 coarse 阶段（3,000 steps）的 mesh 输出。

\textbf{资产生成细节：} 物体 A 使用手机拍摄 28 张多视角照片，经 COLMAP 估计相机位姿，使用 2DGS 训练 4,000 iterations 得到高斯点云重建；物体 B 通过 threestudio 的 dreamfusion-if 配置，以文本 prompt "a cute toy robot, high quality, 3d model" 生成 3D mesh（4,000 iterations）；物体 C 使用手机拍摄单张照片并手工去除背景，通过 Magic123 的 coarse 阶段（3,000 steps）生成 mesh；背景使用 Mip-NeRF 360 的 \texttt{counter} 场景进行 2DGS 重建（4,000 iterations）。

\subsection{统一表示与融合}
不同方法输出的表示包括高斯点、mesh 与纹理 mesh。实验中将资产导出或采样为包含 \texttt{x y z red green blue} 的 ASCII PLY，并通过尺度、旋转和平移参数对齐到同一场景坐标系。融合参数记录于 \texttt{configs/fusion\_transforms.example.json}，包含各物体的 scale、rotation和 translation 向量。

由于各资产生成方法的输出格式差异（2DGS输出高斯点、threestudio输出mesh、Magic123输出mesh），我们统一将资产采样为点云PLY格式进行融合。漫游视频通过 \texttt{render\_fusion\_pointcloud.py} 脚本采用点云投影渲染方式生成（24秒，1280×720分辨率，24fps）。

\subsection{结果与分析}

\begin{table}[H]
\centering
\begin{tabular}{lllll}
\toprule
资产 & 方法 & 顶点/高斯数 & 纹理细节 & 备注 \\
\midrule
物体 A & 多视角 2DGS & 约 150K 高斯 & 高（真实照片） & 28张图像，4000 iter \\
物体 B & 文本到 3D & 约 50K 顶点 & 中（文生纹理） & threestudio DreamFusion \\
物体 C & 单图到 3D & 约 100K 顶点 & 中高（参考图） & Magic123 coarse mesh \\
背景 & 2DGS & 约 300K 高斯 & 高（Mip-NeRF 360） & counter 场景 \\
\midrule
融合场景 & 点云合并 & 615,609 顶点 & 综合 & 统一坐标系对齐 \\
\bottomrule
\end{tabular}
\caption{三种资产生成方式和背景重建对比。}
\label{tab:assets}
\end{table}
  
\textbf{关键帧截图：} 图 \ref{fig:keyframes} 展示了融合场景的6个关键帧，包括视角漫游过程中的不同角度。

\begin{figure}[H]
\centering
\begin{tabular}{ccc}
\includegraphics[width=0.3\textwidth]{fusion_keyframe_01.png} &
\includegraphics[width=0.3\textwidth]{fusion_keyframe_02.png} &
\includegraphics[width=0.3\textwidth]{fusion_keyframe_03.png} \\
(a) 视角1 & (b) 视角2 & (c) 视角3 \\
\includegraphics[width=0.3\textwidth]{fusion_keyframe_04.png} &
\includegraphics[width=0.3\textwidth]{fusion_keyframe_05.png} &
\includegraphics[width=0.3\textwidth]{fusion_keyframe_06.png} \\
(d) 视角4 & (e) 视角5 & (f) 视角6 \\
\end{tabular}
\caption{融合场景漫游视频关键帧（24秒漫游，1280×720分辨率）。}
\label{fig:keyframes}
\end{figure}

\textbf{计算耗时对比：} 表 \ref{tab:cost} 对比了各资产生成方法在计算耗时上的差异。

\begin{table}[H]
\centering
\begin{tabular}{llll}
\toprule
资产 & 方法 & GPU 运行时间 & 备注 \\
\midrule
物体 A & 多视角 2DGS & $\sim$40 min & 28张图像，4000 iter \\
物体 B & threestudio DreamFusion & $\sim$20 min & 4000 iter，SDS Loss \\
物体 C & Magic123 coarse & $\sim$58 min & 3000 coarse steps \\
背景 counter & 2DGS & $\sim$40 min & Mip-NeRF 360 场景 \\
ACT env\_b 训练 & LeRobot ACT & $\sim$4 h & 50K steps，batch 8 \\
ACT env\_abc 训练 & LeRobot ACT & $\sim$5 h & 50K steps，batch 8 \\
\bottomrule
\end{tabular}
\caption{不同方法的 GPU 计算耗时对比（AutoDL V100 单卡）。}
\label{tab:cost}
\end{table}

可以看出，文本生成方法（threestudio）计算开销最小，适合快速原型设计；单图生成方法（Magic123）因需要同时满足 SDS 先验与参考图像约束，耗时较长；真实多视角重建（2DGS）耗时居中。ACT 策略训练因需要处理大量 trajectory 数据并在多帧上做视觉编码，是最耗时的环节（数小时级别）。

\textbf{渲染方式说明：} 由于融合场景包含多种格式的3D资产（高斯点云、mesh），直接使用 Blender 渲染需要复杂的材质转换。本实验采用点云投影渲染方式：将 mesh 采样为点云，与 2DGS 高斯点统一为 PLY 点云格式，使用 Open3D 进行视角变换和投影渲染。这种方式保留了各资产的视觉特征，生成的视频亮度正常（关键帧平均亮度约15-16）。

\section{题目二：LeRobot ACT 跨环境泛化}
\subsection{方法原理}

\textbf{Action Chunking with Transformers (ACT)：} ACT 采用条件变分自编码器（CVAE）架构，以当前观测图像和机器人状态为条件，一次性预测未来 $T=100$ 个时间步的动作序列（即动作分块）。其编码器处理当前时刻和过去若干帧的视觉-状态输入，输出隐变量 $z$；解码器通过 Transformer 架构，以 $z$ 和当前状态为条件，自回归地生成未来 $T$ 步的动作序列。训练损失由两部分组成：动作预测的 L1 损失和 VAE 隐空间的 KL 散度（权重 $\beta=10.0$），即 $\mathcal{L} = \mathcal{L}_{\text{L1}} + \beta \cdot \mathcal{L}_{\text{KL}}$。

\subsection{实验设置}

\begin{table}[H]
\centering
\begin{tabular}{lll}
\toprule
项目 & Baseline (env\_b) & Joint (env\_abc) \\
\midrule
Policy & ACT & ACT \\
Batch size & 8 & 8 \\
Learning rate & 1e-5 & 1e-5 \\
Optimizer & AdamW & AdamW \\
Weight decay & 1e-4 & 1e-4 \\
Training steps & 50,000 & 50,000 \\
Action chunk size & 100 & 100 \\
Hidden dim & 512 & 512 \\
Dim feedforward & 3200 & 3200 \\
Encoder layers & 4 & 4 \\
Decoder layers & 1 & 1 \\
Loss function & L1 + KL ($\beta=10.0$) & L1 + KL ($\beta=10.0$) \\
KL weight & 10.0 & 10.0 \\
Num workers & 4 & 4 \\
Seed & 42 & 42 \\
\bottomrule
\end{tabular}
\caption{ACT 超参数设置。两模型使用完全相同的网络架构和训练超参数，唯一变量为训练数据（splitB vs splitA+B+C）。}
\label{tab:hyperparams}
\end{table}

\subsection{结果}

\begin{table}[H]
\centering
\begin{tabular}{lllll}
\toprule
模型 & 训练环境 & splitD Action L1 $\downarrow$ & splitD Action MSE $\downarrow$ & 测试批次 \\
\midrule
Baseline & B & 0.650 & 0.889 & 381 \\
Joint & A+B+C & \textbf{0.564} & \textbf{0.725} & 254 \\
\bottomrule
\end{tabular}
\caption{环境 D zero-shot 离线评测结果（50 episodes）。}
\label{tab:act_results}
\end{table}

\textbf{Success Rate 说明：} 离线评测 \texttt{eval\_act\_offline.py} 仅能计算 action 误差（L1/MSE），无法直接得到 Success Rate。Success Rate 需要在仿真环境或真实机器人上执行 rollout 并检测任务完成情况才能获取。本实验目前仅提供 action 误差指标，已优于基线（Joint 模型的 Action L1 比 Baseline 低 13.2\%）。

\subsection{Action Chunking 分析}

ACT 的动作分块（Action Chunking）机制通过一次性预测未来 $T$ 个时间步的动作序列（本实验 $T=100$），而非单步预测，具有以下特点：

\textbf{1. 降低短期噪声：} 动作分块通过自回归解码器生成平滑的动作序列，可以有效滤除单步预测的随机噪声。在多环境联合训练中，不同环境的数据分布差异（如物体位置、相机角度）会引入额外的动作噪声，Chunking 机制通过序列建模保持了动作的一致性。

\textbf{2. 视觉误差累积问题：} 长动作块（$T=100$）虽然提高了短期动作连贯性，但也放大了视觉误差累积。当机器人在未见过的环境 D 执行时，前几步的微小视觉估计误差会在后续动作中累积。实验结果显示，联合模型在 splitD 上的 Action L1（0.564）优于基线（0.650），说明多环境训练提高了视觉特征的泛化能力，从而改善了 chunk 内动作的一致性。

\textbf{3. 多环境训练的收益：} 训练于 A+B+C 的联合模型在 zero-shot 测试 splitD 时 action 误差更低，验证了跨环境训练对泛化的正面作用。多环境数据引入了视觉分布偏移（不同环境的光照、物体、背景），使模型学习到更鲁棒的视觉-动作映射，减少了在新环境中的视觉估计误差。

\subsection{训练曲线分析}

\begin{figure}[H]
\centering
\includegraphics[width=0.95\textwidth]{training_loss_total.png}
\caption{ACT 训练 total loss 曲线。两模型均训练 50,000 steps；曲线由本地训练日志每 200 steps 的 \texttt{loss} 字段重建，浅色为原始记录，深色为滑动平均。}
\label{fig:loss_curves}
\end{figure}

\begin{figure}[H]
\centering
\includegraphics[width=0.95\textwidth]{validation_metrics_splitD.png}
\caption{splitD 离线验证指标曲线。上图为 Action L1，下图为 Action MSE；浅色为逐 batch 指标，深色为滑动平均，虚线为整体均值。}
\label{fig:validation_curves}
\end{figure}

从曲线可以观察到：（1）两个模型的 total loss 均快速下降并在后期趋于稳定，说明训练过程正常收敛；（2）Joint 模型在 splitD 上的 Action L1 与 Action MSE 均值低于 Baseline，说明多环境训练提升了 zero-shot 泛化能力；（3）由于未开启在线记录和训练期验证，本实验的验证曲线为最终模型在 splitD 上的离线 batch-wise 指标，而非训练过程中的周期性 validation 曲线。

\section{结论}

\textbf{题目一（2DGS/AIGC融合）：} 完成了多源3D资产生成（真实2DGS、文本生成threestudio、单图生成Magic123）和场景融合，生成了24秒可漫游视频。主要局限包括：仅使用 counter 场景（未包含 garden、bicycle），物体 A 照片数量偏少，以及采用点云投影渲染而非 Blender 真实渲染。

\textbf{题目二（LeRobot ACT）：} 多环境联合训练（A+B+C）在 zero-shot 测试 splitD 上的 action 误差（L1=0.564, MSE=0.725）优于仅使用单环境 B 的基线模型（L1=0.650, MSE=0.889），验证了跨环境训练对策略泛化的积极作用。

\textbf{后续改进：} （1）补充 garden、bicycle 场景的 2DGS 重建；（2）增加物体 A 照片数量并重训；（3）接入仿真环境获取 Success Rate 以补充离线评测指标。

\end{document}
